\SectionStart
\section{问题重述}
\subsection{研究背景与已知条件}
\WritingContract
  {把题目转化为清晰、无方法倾向的研究对象和已知条件。}
  {概括实际背景、对象范围、题目给定的数据或规则，以及会限制答案形式的条件。}
  {以正式题面和附件字段为唯一依据，术语、单位和时间空间范围与题面一致。}
  {不得复制题面长段原文，不得提前介绍算法、结论、实验数字或主观扩展背景。}
\subsection{子问题任务与输出}
\WritingContract
  {让读者明确每一问究竟需要计算、判断、预测或优化什么。}
  {按真实问数逐问说明输入、目标输出、主要约束和交付形式；题面未要求的输出不自行添加。}
  {各问描述与对应 question manifest 的 problem 字段一致。}
  {不得保留不存在的子问题，不得用“建立合理模型”等空话替代可检查的输出。}
\subsection{子问题依赖关系}
\WritingContract
  {说明各问是递进、共享数据还是可以并行求解。}
  {指出前问哪些输出进入后问、共享哪些口径；若无依赖，明确其独立关系。}
  {依赖关系与后文上下游接口及程序数据流一致。}
  {不得在没有题面或模型依据时虚构依赖，不得在本节展开求解流程。}
